%!TEX root = thesis.tex
\chapter{Background and Theory}
\label{chap:background}

% TODO — introductory paragraph

\section{Wind Turbine Aerodynamics and Control}
\label{sec:bg:aero}

\subsection{Power Production and the Betz Limit}
\label{subsec:bg:betz}
% TODO

\subsection{Turbine Control Regions}
\label{subsec:bg:regions}
% TODO

\subsection{Blade Pitch and Generator Torque Control}
\label{subsec:bg:control}
% TODO

\section{SCADA Systems and Condition Monitoring}
\label{sec:bg:scada}

\subsection{What SCADA Systems Measure}
\label{subsec:bg:scada:measure}
% TODO

\subsection{Closed-Loop Data and Its Challenges}
\label{subsec:bg:scada:closedloop}
% TODO

\subsection{Condition Monitoring Using SCADA}
\label{subsec:bg:scada:cm}
% TODO — when explaining that physics-based filtering/preprocessing of SCADA data improves
% model reliability, cite: \citep{hong2025}
% Hong et al. use physics-constrained autoencoders to remove outliers from SCADA power-curve
% data — parallel motivation to the physics-informed zone definitions in Chapter 3.

\subsection{Power Curve Modelling}
\label{subsec:bg:scada:powercurve}
% TODO

\section{Digital Twins in Prognostics and Health Management}
\label{sec:bg:dt}

\subsection{Definitions and Concepts}
\label{subsec:bg:dt:def}
% TODO

\subsection{Requirements for an Interpretable and Adaptive Digital Twin}
\label{subsec:bg:dt:req}
% TODO — when discussing the need to combine physics knowledge with data-driven models
% in a digital twin, cite: \citep{gijon2024}
% Gijón et al. propose a hybrid physics + ML model for wind turbine power prediction that
% achieves 37% accuracy improvement over physics-only while remaining interpretable via SHAP.
% Exact same principle as this thesis: physics sets the structure, data-driven fills the gaps.

\section{Rule-Based Machine Learning}
\label{sec:bg:rules}

\subsection{The RIPPER Algorithm}
\label{subsec:bg:ripper}
% TODO

\subsection{The FOIL Algorithm}
\label{subsec:bg:foil}
% TODO

\subsection{Strengths and Limitations of Rule-Based Methods}
\label{subsec:bg:rules:limits}
% TODO — after explaining the discretization limitation, note that rule-based methods
% remain an active research area. Cite: \citep{nossig2024}
% Nössig et al. (2024) show that RIPPER and FOIL can scale to large datasets using a
% modularity trick (compress → cluster → apply rules per cluster). This shows the algorithms
% are still taken seriously in modern XAI research, not just historical choices.

\section{Supervised Machine Learning Classifiers}
\label{sec:bg:supervised}

\subsection{Random Forests}
\label{subsec:bg:rf}
% TODO

\subsection{Histogram Gradient Boosting}
\label{subsec:bg:gbt}
% TODO

\subsection{Multi-Layer Perceptrons}
\label{subsec:bg:mlp}
% TODO

\section{Deep Reinforcement Learning}
\label{sec:bg:drl}

\subsection{The Reinforcement Learning Framework}
\label{subsec:bg:drl:framework}
% TODO

\subsection{Deep Q-Networks (DQN)}
\label{subsec:bg:dqn}
% TODO

\subsection{Proximal Policy Optimisation (PPO)}
\label{subsec:bg:ppo}
% TODO

\subsection{Advantage Actor-Critic (A2C)}
\label{subsec:bg:a2c}
% TODO

\subsection{DRL in Industrial and Energy Applications}
\label{subsec:bg:drl:industrial}
% TODO — cite BOTH \citep{xie2023} AND \citep{didier2022} here as examples of DRL for
% wind turbine control.
% Xie et al. (2023): DRL (DNN+MPC) for torque and pitch control — already in thesis.
% Didier et al. (2022): DRL (TRPO actor-critic) for pitch control of floating offshore
% wind turbines — shows the same trend applies to the more challenging FOWT case.
% The point: DRL is actively applied to wind energy control; this thesis applies it
% to a different but related problem (classification rather than control).
